\section{Experimental Design}
\label{sec:experiments}

\subsection{Campaign structure and evaluation integrity}

A research-method experimentation campaign iterates through \Evolver{}
revision rounds on a fixed set of development tasks, then evaluates the
frozen final harness against the original initial harness on a separately
registered held-out panel.  Within each registered experiment, foundation
models, provider routes, official evaluators, and Worker execution conditions
are fixed; aligning them creates a comparison condition, not an evolved
capability.

\paragraph{Round structure and promotion.}
Each round runs one \Evolver{} revision followed by a fresh parent and
candidate \Worker{} pair over a declared comparison window.  The \Evolver{}
receives only public task inputs, original Worker artifacts and traces, and
binary outcomes for the development window tasks; it never receives held-out
targets, evaluator code, or reference answers.  The research memory snapshot
advances independently of harness selection: an ABSTAIN or rejected candidate
may still record a methodological lesson.

Candidate promotion uses a Pareto rule over declared benchmark windows: a
complete candidate is retained only when no benchmark aggregate key regresses
and at least one strictly improves; exact ties and mixed gains/regressions
retain the earlier whole parent.  This is not a per-task non-regression
guarantee.  No task-wise or component-wise best-version splicing occurs.

\paragraph{Promotion gate calibration.}
Single-draw comparisons have high false-promotion rates under a zero-threshold
gate.  On a panel where 68\% of tasks are deterministic and the mean flippable
task count per 12-task window is 3.5, a gate requiring any positive mean
difference promotes a neutral mechanism roughly 13\% of the time per round,
yielding about 1.6 spurious promotions per six-round campaign.  A net-\(\geq
2\) gains rule on the target window reduces this to 0.6 at identical cost, and
is the target gate for future campaign registrations (Section~\ref{sec:limitations}).
The current BacktestBench comparison uses a whole-package selection rule;
the power properties of that specific gate are reported alongside the result.

\paragraph{Initial-versus-frozen endpoint.}
After development, exactly one complete \(\Hfinal\) is frozen and evaluated
with fresh \Worker{} instances against the registered \(\Hinit\) on a
separately fixed panel that was not consulted during development.  Historical
score vectors from development rounds cannot substitute for this fresh
measurement.  The completed BacktestBench endpoint (Section~\ref{sec:current-results})
uses 20 released test instances fixed before any test outcome is observed,
with the actual complete initial harness \texttt{fe81e756} versus the unchanged
promoted candidate \texttt{403fa470}.

\paragraph{Test-time-scaling baseline.}
Harness evolution should be compared against best-of-\(K\) parallel sampling
from the frozen initial harness at matched total Worker compute
\citep{rethinking2026harness}.  The current campaigns do not include this
arm.  It is a registered requirement for the next campaign and will be
reported when complete.

\subsection{Benchmarks and task selection}

QFBench contains practitioner-style quantitative assignments
\citep{qfbench2026}; QuantCodeEval evaluates strategy-code reproduction from
finance papers and task instructions \citep{quantcodeeval2026}.  They provide
operational tests of quantitative work.  Success does not establish investment
profitability, original alpha discovery, or unrestricted scientific expertise.

\paragraph{BacktestBench.}
BacktestBench supplies 18{,}246 questions in four categories from over six
million market records \citep{wang2026backtestbench}.  Our screen uses
12 development tasks (eight train, four validation) and a 20-task frozen
endpoint panel fixed before test observation.  The public selector fixes
20 released test tasks (five per category) by lexicographic UUID order,
excluding all 32 previously declared instances and normalized public-text
clones.  Input-availability checks establish usable public market rows;
they do not prove disjoint templates, market windows, or latent lineage.

The \Worker{} receives the raw market CSV archive and public catalogs rather
than the published SQL/factor-retrieval AutoBacktest pipeline.  Comparison
operators, a declared stock-selection category alias, and file-first answer
precedence are preserved.  The public source is pinned to revision
\texttt{4bacdfbe83b3}.  Revision evidence includes only the eight train tasks,
their original \Worker{} artifacts and traces, and binary outcomes; the
complete inherited harness remains mutable.

\paragraph{Selection rule for the BacktestBench endpoint.}
Eligible whole admitted candidate packages are ranked on the fixed
development twelve by known successes, then fewer unresolved cells, then
fewer known regressions, then earlier registration.  Null is neither a
success nor a regression.  Exactly one complete package is frozen before
any of the 20 test outcomes is read.  The rule selects the cap-only E
(frozen before test entry); the initial versus frozen comparison closes with
all 40 cells scored.

\paragraph{QuantCodeEval and QFBench.}
The registered mixed campaign evolves one whole shared harness across both
benchmarks.  Selection uses each benchmark's declared native aggregate
without pooling: QCE ranks by binary-success count then equal-task mean
fraction; QFBench ranks by equal-task mean official reward then fraction.
The QCE/QF operation-family condition uses T24/T26/T27/T29 and 13F; missing
original measurements remain null.

\subsection{Evolver revision policy}

The \Evolver{} may edit any of the nine component roles in the candidate
harness: system prompt, agent configuration, tool descriptions, tool
implementations, validator, skills, memory, middleware, and routing.  All
nine are equally legitimate mutation targets; the revision is not restricted
to prompts.

For computational failure classes --- \texttt{temporal\_causality},
\texttt{portfolio\_accounting}, \texttt{formula\_parameterization}, and
\texttt{data\_universe\_preprocessing} --- the revision policy now requires
at least one executable component (\texttt{tools/}, \texttt{validator/}, or
\texttt{skills/}) alongside any prompt text.  A prompt-only edit for these
classes does not establish a mechanism test, because a Worker following the
edited prompt can still produce the same failure on a different task instance
with different numbers or windows.  This constraint was added on October 6
2026 in response to the observation that the BacktestBench positive result
came from a prompt-only revision.  It targets future rounds, not the
completed endpoint.

\subsection{Evidence scope and comparison conditions}

Within a comparison, initial and frozen harnesses receive matched turn
budget, context and output limits, provider routes, timeouts, and runtime
allowances.  Raising or aligning an inference allowance for both arms is a
new comparison condition, not an evolved capability.

The October 6 cap-only condition permits declared training scalar targets to
the \Evolver{} after blind development attempts; its \Worker{}s remain
label-free.  The default and every earlier public-only condition remain
answer-free.  Both prospective arms receive the same public output-consumer
clarification; an interface-only benefit is not a numerical research gain.

All valid outcomes, failed attempts, and costs remain visible.  Adaptive
development selection can overfit evaluation criteria
\citep{cawley2010selectionbias}; the separately registered initial-versus-frozen
endpoint is the primary performance evidence and is not adjusted by
development observations.
